\section{Spatial isometries and nearby Kerr centers}
\label{phase:symmetry}

Return to the arbitrary fixed center $M>0$, $0<\mathfrak a<M$ and
its datum $d_*=(h_*,K_*)$ on $\Sigma$. In this appendix $\mathcal D$
is the corresponding weighted constraint space of
Definition~\ref{phase:space}, before the mass normalization.

The neighborhood is nonsymmetric in the sense that no symmetry is
imposed on its data; it necessarily contains the symmetric center.
It also contains data with trivial spatial isometry group. We give
the elementary construction to distinguish this assertion from any
claim about spacetime Killing fields.

\begin{lemma}
\label{phase:trivial-isometries}
Every neighborhood of \(d_*\) in \(\mathcal D\) contains complete
vacuum data whose spatial data isometry group is trivial.
\end{lemma}

\begin{proof}
The reference slice is maximal: its second-form trace vanishes.
For example, simultaneous reversal of Killing time and azimuth
preserves the spatial data convention while reversing this scalar,
which is axisymmetric. A compactly supported small spacelike
normal graph in the smooth Kerr spacetime gives exact vacuum
constraint data, identical to \(d_*\) off a compact set. The first
variation of its second-form trace is an operator with principal
part \(\pm\Delta_{h_*}\), plus smooth lower-order terms.

Choose a small normal-coordinate ball and graph height
\(f_\ell(z)=\ell^2f_0(z/\ell)\).
After rescaling to the fixed coordinate ball, the linearized trace
converges smoothly, as \(\ell\to0\), to
\(\pm\Delta_{\R^3}f_0\).
One can choose \(f_0\) compactly supported so that this last
function has a unique strict positive global maximum, with
nondegenerate Hessian of simple spectrum, and with nonzero third
derivative in each of the three Hessian eigendirections.
Here is a direct choice. Start with a negative Gaussian whose three
axis scales are distinct and sufficiently close; its Laplacian has
a unique strict maximum at the origin and a negative Hessian with
simple eigenvalues. Cut it off sufficiently far away, where all
its derivatives are small, and add arbitrarily small cutoff fifth
powers in the coordinate directions. The Laplacian of each fifth
power has a nonzero cubic derivative in its own direction. The
strict maximum and simple Hessian persist.

First take \(\ell\) small. Then take the actual normal-graph
amplitude small enough that its trace, divided by that amplitude,
has the same properties. All statements concern the induced
metric, whose rescaled coefficients and derivatives converge to
the Euclidean ones on the patch. Outside the patch the trace
is zero, below the selected positive maximum.

A data isometry must fix the unique maximum. Its differential
preserves the Hessian, hence each one-dimensional eigenspace.
The nonzero odd third derivatives exclude each possible sign
flip, so its differential at the fixed point is the identity.
An isometry of a connected Riemannian manifold is determined by
its value and differential at one point, and is therefore the
identity. The graph is unchanged outside a compact set and remains
complete. Making its amplitude tend to zero gives convergence in
every prescribed finite set of seminorms.
\end{proof}

These examples are hypersurfaces in Kerr and do not remove its
spacetime Killing fields. The perturbations used to prove
Theorem~\ref{phase:main} are the separate signed construction of
Sections~\ref{prep:section}--\ref{holo:section}.

Nearby Kerr parameters are also allowed. Identify the bridges near
the waist using a signed square root of \(r-r_+(M,\mathfrak a)\).
The pole in the Boyer--Lindquist \(dr^2\) coefficient then becomes
a smooth positive radial metric coefficient. The angular velocity
of the azimuthal shift is smooth in \(r,\theta\) and constant on the
horizon sphere; dividing its metric Lie derivative by the signed
lapse gives a smooth second form across the waist. The end formulas
have parameter-smooth symbol estimates with fixed Euclidean leading
metric. Thus these identifications place nearby Kerr centers in
the same phase space.
